\chapter{Taxonomy Generation Engine}
\label{chap:taxonomy}

\section{Algorithm 1: Reasoning-Driven Taxonomy Generation}

The taxonomy builder implements Algorithm~1 from the Simula research paper. Given a
dataset description $y$ and target depth $D$, it produces a set of taxonomies
$\mathcal{T} = \{T_i\}_{i=0}^{K}$ that map the conceptual space.

\begin{figure}[H]
\centering
\begin{tikzpicture}[
    node distance=0.6cm and 1.5cm,
    stepbox/.style={rectangle, draw, rounded corners, fill=llmpurple!10,
                 minimum width=3cm, minimum height=0.7cm, font=\small,
                 text width=2.8cm, align=center},
    arr/.style={-{Stealth[length=2mm]}, thick},
]
    \node[stepbox] (f1) {1. Factor Identification};
    \node[stepbox, below=of f1] (f2) {2. Breadth-First Expansion};
    \node[stepbox, below=of f2] (f3) {3. Best-of-N Proposal};
    \node[stepbox, below=of f3] (f4) {4. Critic Refinement};
    \node[stepbox, below=of f4] (f5) {5. Level Planning};
    
    \draw[arr] (f1) -- (f2);
    \draw[arr] (f2) -- (f3);
    \draw[arr] (f3) -- (f4);
    \draw[arr] (f4) -- (f5);
    \draw[arr, dashed] (f5.east) -- ++(1,0) |- node[right, pos=0.25, font=\tiny]{next level} (f2.east);
\end{tikzpicture}
\caption{Taxonomy generation stages with iterative refinement.}
\label{fig:taxonomy-stages}
\end{figure}

\section{Step 1: Factor Identification}

Given schema metadata, the LLM proposes factors of variation:

\begin{lstlisting}[language=jsonx,caption={Factor identification prompt}]
SYSTEM: You identify the prime factors of variation for a 
        Text-to-SQL training dataset.
        
USER:   Schema context:
        {schema_context}
        
        Identify 3-7 factors that capture the main dimensions
        of queries users might ask. Examples:
        - "query_complexity" (simple, join, aggregation, subquery)
        - "time_range" (point-in-time, period, YTD, MTD)
        - "entity_type" (customer, order, product, supplier)
        
        Output JSON: {"factors": [{"name": "...", "description": "..."}]}
\end{lstlisting}

\section{Step 2: Breadth-First Expansion}

Each factor is expanded into a taxonomy tree. For each node, the LLM proposes children:

\begin{lstlisting}[language=jsonx,caption={Node expansion (simula\_taxonomy\_builder.py, lines 180--220)}]
async def _expand_node(
    self, node: TaxonomyNode, context: str
) -> list[TaxonomyNode]:
    """Expand a single node with Best-of-N proposals."""
    
    # Step 1: Generate N proposals
    proposals = await self.llm.best_of_n(
        prompt=self._build_expansion_prompt(node, context),
        n=self.config.taxonomy.best_of_n,
        temperature=0.8,
    )
    
    # Merge proposals into candidate set
    candidates = self._merge_proposals(proposals)
    
    # Step 2: Critic refinement
    refined = await self._critic_refine(node, candidates, context)
    
    return refined
\end{lstlisting}

\section{Step 3: Best-of-N Proposal}

Multiple proposals are generated to increase coverage:

\begin{lstlisting}[language=jsonx,caption={Best-of-N sampling}]
async def best_of_n(
    self,
    prompt: str,
    n: int = 5,
    temperature: float = 0.8,
) -> list[str]:
    """Generate N diverse completions."""
    tasks = [
        self.generate(prompt, temperature=temperature)
        for _ in range(n)
    ]
    responses = await asyncio.gather(*tasks)
    return [r.content for r in responses]
\end{lstlisting}

\textbf{Known limitation:} All $N$ proposals use the same temperature, which may reduce 
diversity compared to varying temperature across samples. Empirical testing suggests 
this is acceptable for taxonomy generation where structural diversity (prompt variation) 
dominates over sampling diversity. v2 will evaluate temperature annealing ($T \in [0.6, 1.0]$).

\section{Step 4: Critic Refinement}

The critic evaluates and refines the proposed children:

\begin{lstlisting}[language=jsonx,caption={Critic refinement prompt}]
SYSTEM: You are a taxonomy critic. Evaluate the proposed children
        for COMPLETENESS, SOUNDNESS, and SPECIFICITY.
        
USER:   Parent node: {parent.name} - {parent.description}
        Siblings: {siblings}
        Proposed children: {candidates}
        
        Actions you may take:
        - ADD: missing concepts
        - REMOVE: irrelevant or duplicate nodes
        - MERGE: overlapping concepts
        - EDIT: improve descriptions
        
        Output JSON with refined children list.
\end{lstlisting}

\section{Step 5: Level Planning}

After completing a level, a plan is generated for the next level:

\begin{lstlisting}[language=jsonx,caption={Level planning}]
async def _plan_next_level(
    self, current_level_nodes: list[TaxonomyNode]
) -> str:
    """Generate guidance for next level expansion."""
    prompt = f"""
    Current level nodes: {[n.name for n in current_level_nodes]}
    
    Generate a plan for expanding the next level that ensures:
    1. Consistent granularity across all branches
    2. No overlap between sibling expansions
    3. Complete coverage of the concept space
    """
    response = await self.llm.generate(prompt)
    return response.content
\end{lstlisting}

\section{Taxonomy Depth and Dataset Size Relationship}
\label{sec:depth-size-relationship}

The relationship between taxonomy depth ($D$), number of factors ($K$), and target
dataset size ($N$) is critical for planning generation runs. This section provides
guidance for configuring these parameters.

\subsection{Node Count Growth}

Taxonomy nodes grow exponentially with depth. Given an average branching factor $b$
(typically 3--5 for well-structured taxonomies), the total node count at depth $D$ is:

\[
\text{Total Nodes} = K \times \sum_{d=0}^{D} b^d = K \times \frac{b^{D+1} - 1}{b - 1}
\]

\noindent
For $K=4$ factors and $b=4$ average branching:

\begin{table}[H]
\centering
\caption{Expected Node Count by Taxonomy Depth}
\label{tab:depth-nodes}
\begin{tabular}{llll}
\toprule
\textbf{Depth} & \textbf{Nodes/Factor} & \textbf{Total Nodes (K=4)} & \textbf{Leaf Nodes} \\
\midrule
1 & 5 & 20 & 16 \\
2 & 21 & 84 & 64 \\
3 & 85 & 340 & 256 \\
4 & 341 & 1,364 & 1,024 \\
5 & 1,365 & 5,460 & 4,096 \\
\bottomrule
\end{tabular}
\end{table}

\subsection{Target Examples vs Taxonomy Depth}

To achieve adequate coverage, each leaf node should have multiple examples. The
recommended formula for minimum examples:

\[
N_{\min} = L \times m \times (1 + r)
\]

\noindent
Where:
\begin{itemize}
  \item $L$ = number of leaf nodes (approximately $K \times b^D$)
  \item $m$ = minimum examples per unique mix (recommended: 5--10)
  \item $r$ = critic rejection rate (typically 0.08--0.15)
\end{itemize}

\begin{table}[H]
\centering
\caption{Recommended Target Examples by Depth (K=4, b=4, m=5, r=0.10)}
\label{tab:target-by-depth}
\begin{tabular}{lllll}
\toprule
\textbf{Depth} & \textbf{Leaf Nodes} & \textbf{Unique Mixes} & \textbf{Min Examples} & \textbf{Recommended} \\
\midrule
2 & 64 & ~100 & 550 & 1,000 \\
3 & 256 & ~400 & 2,200 & 5,000 \\
4 & 1,024 & ~1,600 & 8,800 & 10,000--15,000 \\
5 & 4,096 & ~6,500 & 35,750 & 50,000+ \\
\bottomrule
\end{tabular}
\end{table}

\subsection{Configuration Guidelines}

\begin{description}
  \item[Small pilot (1,000 examples):] Use depth 2, 3--4 factors. Focus on 
        demonstrating coverage, not exhaustive diversity.
  \item[Standard production (10,000 examples):] Use depth 3--4, 4--5 factors.
        Balances coverage with generation cost.
  \item[Large-scale training (50,000+ examples):] Use depth 4--5, 5--7 factors.
        Requires significant LLM compute budget.
\end{description}

\subsection{Cost Estimation}

Each generated example requires approximately:
\begin{itemize}
  \item 1--2 LLM calls for meta-prompt generation
  \item 1 LLM call for example generation
  \item 2 LLM calls for double-critic evaluation
  \item 0--1 LLM calls for complexification (50\% of examples)
\end{itemize}

\noindent
Total: approximately 4.5 LLM calls per accepted example (accounting for rejections).
At \$0.002/1K tokens and ~1K tokens per call, the cost is approximately
\$0.01 per accepted example, or \$100 per 10,000 examples.

\section{Worked Example: Trial Balance Taxonomy}

\begin{table}[H]
\centering
\caption{Sample Taxonomy for Trial Balance Domain}
\label{tab:tb-taxonomy}
\begin{tabular}{llll}
\toprule
\textbf{Level 0} & \textbf{Level 1} & \textbf{Level 2} & \textbf{Level 3} \\
\midrule
query\_type & aggregation & sum & by\_account \\
            &             &     & by\_department \\
            &             & count & distinct\_accounts \\
            & comparison & variance & MoM \\
            &            &          & YoY \\
            & filter & by\_threshold & exceeds\_100M \\
\bottomrule
\end{tabular}
\end{table}

% =============================================================================
% CHAPTER 5: TRAINING DATA GENERATION ENGINE
% =============================================================================
\newpage